%%%%%%%%%%%%%%%%%%%%%%%%%%%%%%
%%%%%%%%%%%% Preamble %%%%%%%%%%%%%
%%%%%%%%%%%%%%%%%%%%%%%%%%%%%%

% Declare document class and miscellaneous packages
\documentclass[english]{article}
\usepackage{natbib}
\usepackage[leqno]{amsmath}
\usepackage{amssymb}
\usepackage{mathrsfs}
\usepackage{mathtools}
\usepackage{setspace}
\usepackage{fullpage}
\usepackage[margin=1in]{geometry}
\interfootnotelinepenalty=10000

%Hyper-references
\usepackage{hyperref}
\hypersetup{colorlinks, citecolor=black, filecolor=black, linkcolor=black, 
urlcolor=black, pdftex}

\begin{document}

\doublespacing

%%%%%%%%%%%%%%%%%%%%%%%%%%%%%%
%%%%%%%%%%%%% Body %%%%%%%%%%%%%%
%%%%%%%%%%%%%%%%%%%%%%%%%%%%%%

Originally developed by \citet{friedman2001greedy}, gradient boosting machines are a family of 
machine-learning techniques that combine a large number of weak learners to form 
a stronger ensemble prediction.
Unlike some common ensemble techniques (e.g., random forests) that simply average models
in the ensemble, gradient boosting adds models to the ensemble sequentially by fitting new
weak learners to the negative gradient of some chosen loss function.
With the classic squared-error loss function, this amounts to sequentially fitting new models to 
the current residuals of the ensemble.%
\footnote{See \citet{natekin2013gradient} for an accessible introduction to gradient boosting.}
Extreme gradient boosting (EGB) is a particularly popular implementation of gradient boosting
that uses classification and regression trees as the base models or weak learners.
The popularity of EGB is due to the fact that it is fast, scalable, and has been shown to outperform 
competing methods across a wide range of problems \citep{chen2016xgboost}.

Classification and regression trees are based on a sequential binary splitting of the covariate space that
serves to partition the data into subsamples that are increasingly homogeneous in terms of the outcome 
variable.
A tree begins with a single root node that is split into two child nodes according to some decision
rule.
A decision rule pertains to a single explanatory variable and observations are assigned to the 
child nodes depending on whether they meet some condition associated with the explanatory
variable.%
\footnote{For example, for a continuous explanatory variable, one determines whether a given
observations falls above or below some numeric cutoff.} 
The child nodes can be further split into new child nodes based on different decision rules that
involve other explanatory variables and cutoff points. 
Any child node that is not further split is referred to as a terminal node or leaf, and each leaf is
associated with a parameter that provides a prediction for the subsample associated
with that leaf.
For any given observation, the ensemble prediction used by EGB is the sum of that observation's
predictions across all trees in the ensemble.

In what follows, we discuss the key features of the EGB algorithm, drawing on the presentation provided 
in \citet{chen2016xgboost}.
Let $y_i$ denote the outcome of interest (e.g., poverty headcounts) for observation or entity $i$, and let 
$x_i$ denote a vector of explanatory variables.
In line with the above, the predicted value of the outcome $\hat{y}_i$ is given by the sum of predictions
across the individual trees in the ensemble:
\begin{equation}
\hat{y}_i = \sum_t f_t(x_i)
\end{equation}
where the trees are indexed by $t$ and $f_t(x_i)$ gives the prediction of a given tree as a function of 
$x_i$.
To build the trees used in the model, EGB minimizes the following regularized objective function:
\begin{equation}
\mathcal{L} = \sum_i l(y_i, \hat{y}_i) + \sum_t r(f_t)
\end{equation}
where $l(y_i, \hat{y}_i)$ is the loss associated with the $i^{th}$ observation and $r(f_t)$ is a 
regularization term that serves to mitigate overfitting.
In our implementation, we use the classic squared-error loss such that $l(y_i, \hat{y}_i) = 
(y_i - \hat{y}_i)^2$.

EGB uses the following specification of the regularization term:
\begin{equation}
r(f_t) = \gamma M_t +  \frac{1}{2}\lambda \sum_j m_{tj}^2
\end{equation}
where $M_t$ is the total number of leaves in tree $t$, $m_{tj}$ is the prediction
associated with the $j^{th}$ leaf in a given tree, and $\gamma$ and $\lambda$ are 
hyperparameters that must be chosen by the user (we discuss our selection of
hyperparameters below).
The regularization term serves to mitigate overfitting by penalizing complicated trees 
and smoothing the predictions associated with each tree.
That is, the method relies on a large number of weak learners and the regularization 
term helps control the complexity of the learners.
It is important to note that the above specification of the regularization term is only one of many 
possible specifications, though \citet{chen2016xgboost} claim that it is simpler than the
leading alternatives and easier to parallelize.
Further note that when the hyperparameters are set to zero, the objective function becomes
equivalent to that used in traditional gradient tree boosting.

To minimize the objective function, the model is trained in a sequential manner, building one
tree at a time. 
The objective function at iteration $t$ can be written as follows:
\begin{equation}
\mathcal{L}_t = \sum_i l(y_i, \hat{y}_i^{t-1} + f_t(x_i)) + r(f_t)
\end{equation}
where $\hat{y}_i^t = \hat{y}_i^{t-1} + f_t(x_i)$.
Each iteration then adds the tree $f_t(x_i)$ that greedily minimizes the objective function at that
iteration.
More specifically, EGB uses a second-order Taylor series expansion to approximate the loss
function at any given iteration:
\begin{equation}
\tilde{\mathcal{L}}_t = \sum_i \left\{g(y_i, \hat{y}_i^{t-1}) f_t(x_i) 
+ \frac{1}{2}h(y_i, \hat{y}_i^{t-1}) f_t(x_i)^2 \right\} + r(f_t)
\label{eq:approx}
\end{equation}
where $g(\cdot, \cdot)$ and $h(\cdot, \cdot)$ denote the first- and second-order derivatives of the 
loss function $l(y_i, \hat{y}_i^{t-1})$ with respect to $\hat{y}_i^{t-1}$.
That is, the point $\hat{y}_i^{t-1}$ is taken as the basis for the expansion.
The above conveniently omits the first term of the series $l(y_i, \hat{y}_i^{t-1})$ because it is a constant.

The EGB algorithm is based on a re-expression of Eq.\ (\ref{eq:approx}).
Let $s_j$ denote the set of observations that belong to the $j^{th}$ leaf of the tree being built at 
iteration $t$.
Substituting in the explicit form of the regularization term and noting that all observation that belong to 
the same leaf receive the same prediction $m_{tj}$, we can write
\begin{equation}
\tilde{\mathcal{L}}_t = \sum_j \sum_{i \in s_j} \left\{g(y_i, \hat{y}_i^{t-1}) m_{tj}
+ \frac{1}{2}h(y_i, \hat{y}_i^{t-1}) m_{tj}^2 \right\} + \gamma M_t +  \frac{1}{2}\lambda \sum_j m_{tj}^2
\end{equation}
or, equivalently,
\begin{equation}
\tilde{\mathcal{L}}_t = \sum_j \left\{G_j m_{tj} + \frac{1}{2}(H_j + \lambda)m_{tj}^2  \right\} + \gamma M_t
\end{equation}
where $G_j$ and $H_j$ represent the sums over $g(\cdot, \cdot)$ and $h(\cdot, \cdot)$ for all observations
in $s_j$.
Note that the summation in the above is over $j$, which indexes the leaves in the tree.

For a given tree structure, we can find the optimal prediction values by minimizing the objective function 
with respect to $m_{tj}$, which yields $m_{tj}^* = - G_j/(H_j + \lambda)$.
With a squared-error loss function, where $g(y_i, \hat{y}_i^{t-1}) = -2(y_i - \hat{y}_i^{t-1})$ and 
$h(y_i, \hat{y}_i^{t-1}) = 2$, the optimal prediction values are effectively the (penalized) average of the 
residuals associated with a given leaf.
We can then substitute the optimal prediction values into the objective function to obtain the following: 
\begin{equation}
\tilde{\mathcal{L}}_t^* = - \frac{1}{2}\sum_j \frac{G_j^2}{H_j + \lambda} + \gamma M_t
\end{equation}
which gives us the minimal objective function for a given tree structure.
Ideally, one would enumerate all possible tree structures and then use the tree associated with 
the smallest value of the (minimized) objective function, but this is often intractable.
EGB instead implements an alternative approach that seeks to minimize the objective at a given iteration 
by greedily optimizing one level of the tree at a time.

To this end, consider splitting a node in a given tree into ``left'' and ``right'' child nodes based on some
candidate splitting rule.
Let $G_L$ and $G_R$ represent $G_j$ for the left and right child nodes, respectively.
Similarly, let $H_L$ and $H_R$ represent $H_j$ for the child nodes.
We can then find the reduction in the objective function as follows:
\begin{equation}
\Delta \tilde{\mathcal{L}}_t^* = \frac{1}{2}\left\{\frac{G_L^2}{H_L + \lambda}  + \frac{G_R^2}{H_R + \lambda}
- \frac{(G_L + G_R)^2}{H_L + H_R + \lambda} \right\} - \gamma
\label{eq:gain}
\end{equation}
where the first two terms in parentheses capture the loss associated with the child nodes and the third
term captures the loss associated with the parent node being split.
When splitting a node in any given tree, EGB then evaluates Eq.\ (\ref{eq:gain}) for every possible split
rule for every available explanatory variable and chooses the split associated with the maximal reduction in 
the objective function.%
\footnote{This is known as the ``exact greedy algorithm.'' 
EGB also has an approximate algorithm that can be used for large datasets.
See \citet{chen2016xgboost} for details.}
EGB then continues to split nodes in a given tree until a stopping criterion is met (e.g., a user-specified
maximum tree depth).
Once a given tree is built, EGB then moves to building the next tree in the ensemble and continues to do
so until reaching the user-specified number of trees to build.

The above references several hyperparameters that must be chosen by the user, including the 
hyperparameters in the regularization term (i.e., $\gamma$ and $\lambda$), the maximum tree depth, and 
the number of trees in the ensemble. 
There are two additional hyperparameters that can be used to further prevent overfitting.
First, one can use column or feature subsampling where each tree in the ensemble is built based on a
(uniformly) random subset of all explanatory variables.
This procedure can also speed up the algorithm because it reduces the number of covariates that 
must be searched across when creating new splits.
Second, one can use shrinkage by scaling the predictions associated with each tree by a factor $\eta$,
which is often called the learning rate.
The predicted value of the outcome at iteration $t$ is then given by $\hat{y}_i^t = 
\hat{y}_i^{t-1} + \eta f_t(x_i)$.
Shrinkage reduces the influence of each tree so that no individual tree has too much influence on the
ensemble prediction.

One possible approach to hyperparameter selection is to use a grid search where the user (1) specifies
the domain of the hyperparameters, (2) applies the model to every combination of hyperparameters,
and (3) selects the hyperparameters that perform the best in terms of some metric (e.g., mean squared
prediction error in the validation dataset).
Grid search, however, is computationally inefficient because hyperparameter proposals are uninformed
by previous evaluations, meaning that the procedure spends considerable time evaluating suboptimal
hyperparameters.
We instead use a Bayesian hyperparameter optimization procedure, which sequentially updates a probability 
model that relates the hyperparameter space to the performance metric and then uses the model to 
intelligently explore the hyperparameter space \citep{bergstra2011algorithms}.
Specifically, our implementation of gradient boosting relies on the open-source library XGBoost, which
we combine with the Hyperopt library that conducts Bayesian hyperparameter optimization using
the tree-structured Parzen estimator.

% Notation
% g, G: Gradient
% gamma: Hyperparameter on the number of leaves
% h, H: Hessian
% i: Indexes observations
% j: Indexes for the leaves
% l: Loss function
% lambda: Hyperparameters on the leaf weights
% m: Prediction values
% r: Regularization function
% s: Set of observations in a given leaf
% t: Indexes trees
% x: Explanatory variables
% y: Outcome


%%%%%%%%%%%%%%%%%%%%%%%%%%%%%%
%%%%%%%%%%%% References %%%%%%%%%%%%
%%%%%%%%%%%%%%%%%%%%%%%%%%%%%%

%Start fresh page
\newpage
\cleardoublepage

%Declare the style of the bibliography
\bibliographystyle{apalike}
\bibliography{References}

\end{document}